\section{Related Work}
\label{sec:related}

\subsection{Secure Neural-Network Inference}
Gazelle combined homomorphic encryption for linear layers with garbled circuits for nonlinear layers, establishing a practical hybrid design for two-party neural-network inference \cite{juvekar2018gazelle}. Delphi paired secure protocols with an architecture-search component that selects network structures according to secure-inference cost \cite{mishra2020delphi}. ABY3 developed a mixed-protocol framework for three-party machine learning, supporting conversions among arithmetic, Boolean, and binary-circuit representations \cite{mohassel2018aby3}. These systems established core principles that remain relevant to Transformer inference: linear and nonlinear operators should not be forced into a single cryptographic representation, and conversion costs must be included in system design.

Our setting differs in its focus on a high-throughput, semi-honest 3PC implementation for Transformer encoder inference. We build on hpmpc, which studies high-throughput honest-majority MPC across network settings, and PIGEON, which adds model import, neural-network layers, and protocol-independent inference to hpmpc \cite{harth2025hpmpc,harth2025pigeon}. Rather than proposing a new general-purpose sharing scheme, this work adds Transformer-oriented operators and an accelerator-aware execution path to the existing stack.

\subsection{Private Transformer Inference}
Iron was among the first systems to address private Transformer inference directly. It introduced protocols for matrix multiplication and the nonlinear Softmax, GELU, and LayerNorm functions in a two-party client--server design \cite{hao2022iron}. MPCFormer replaced costly nonlinearities with MPC-friendly approximations and used knowledge distillation to reduce the resulting quality loss \cite{li2023mpcformer}. PUMA developed secure protocols for Transformer functions, including embedding, Softmax, GELU, and LayerNorm, and emphasized inference without task-specific model retraining \cite{dong2023puma}. Ditto pursued a complementary direction by combining quantization-aware distillation, efficient representation conversion, and secure quantized operators \cite{wu2024ditto}.

These systems expose three broad design choices. First, a protocol may preserve the pretrained model and invest more work in high-quality secure approximations. Second, it may modify or retrain the model to reduce cryptographic cost. Third, it may reduce precision through quantization and design protocols around the resulting representation. Our design follows the first path for the initial BERT-base target: it retains the standard attention structure, uses a 32-bit fixed-point representation, and approximates nonlinear functions inside the secure protocol. This choice isolates systems and numerical effects from retraining effects, although future quantized variants remain possible.

\subsection{GPU-Accelerated MPC}
CryptGPU executes privacy-preserving machine-learning operations over secret-shared values with CUDA kernels and demonstrates the value of keeping secure tensors on the GPU \cite{tan2021cryptgpu}. Piranha provides a modular GPU platform for secret-sharing protocols and supports several party counts and security models through reusable protocol and kernel layers \cite{watson2022piranha}. PIGEON incorporates CUDA matrix multiplication and convolution into a high-throughput inference framework \cite{harth2025pigeon}. Our work is closest to this line of systems research, but its unit of optimization is a Transformer block: GPU matrix multiplication must be coordinated with comparison-heavy Softmax, normalization, truncation, and residual paths rather than evaluated as an isolated accelerator kernel.

\subsection{Positioning}
The proposed system occupies the intersection of three lines of work: mixed-protocol 3PC, private Transformer operators, and GPU-oriented secure computation. Relative to prior private-Transformer systems, the intended contribution is an implementation and evaluation in the hpmpc/PIGEON/Trio stack with explicit accelerator and network measurements. Relative to general GPU-MPC systems, the intended contribution is a Transformer-specific operator pipeline. The empirical magnitude of any latency, throughput, communication, or utilization improvement is deliberately left open until the full BERT-base pipeline and controlled baselines are available.
